\subsection{包络代数的中心}

在本节中，我们选取 R 的一个基 B。令 \(R_{+}\) 为关于 B 的正根集合。令 \(\rho = \frac{1}{2} \sum_{\alpha \in R_{+}} \alpha\) ，并令 \(\delta\) 为代数 \(\mathbf{S}(\mathfrak{h})\) 的自同构，它将每个 \(x \in h\) 变为 \(x - \rho(x)\) ，从而将 \(h^{*}\) 上的多项式函数 p 变为函数 \(\lambda \mapsto p(\lambda - \rho)\) 。

定理 2。令 U 为 g 的包络代数，Z 为其中心，V ⊂ U 为 h 的包络代数（视为 S(h)），U⁰ 为 V 在 U 中的交换子，\(\varphi\) 为相对于 B 从 \(U^{0}\) 到 V 的 Harish-Chandra 同态（§6，第4号）。令 \(\mathbf{S}(\mathfrak{h})^{\mathrm{W}}\) 为 \(\mathbf{S}(\mathfrak{h})\) 中在 W 的作用下不变的元素所成的集合。则 \((\delta\circ\varphi)|Z\) 是从 Z 到 \(\mathbf{S}(\mathfrak{h})^{\mathrm{W}}\) 的同构，且与 B 的选择无关。

\begin{enumerate}
\def\labelenumi{\alph{enumi})}
\tightlist
\item
  令 \(P_{++}\) 为 R 的支配权集合，\(w \in W\)，\(\lambda \in P_{++}\)，\(\mu = w\lambda\)。则 \(Z(\mu - \rho)\) 同构于 \(Z(\lambda - \rho)\) 的一个子模（§6，第3号，命题6的推论2），并且对所有 \(\varphi(u)(\lambda - \rho) = \varphi(u)(\mu - \rho)\) 都有 \(u \in Z\)（§6，第4号，命题7）。因此，\((\delta \circ \varphi)(u)\) 上的多项式函数 \((\delta \circ \varphi)(u) \circ w\) 与 \(h^{*}\) 在 \(P_{++}\) 上相同。但 \(P_{++}\) 在 Zariski 拓扑下稠密于 \(h^{*}\)：这是因为可以利用由基本权 \(h^{*}\) 组成的基将 \(k^{B}\) 识别为 \(\varpi_{\alpha}\)，再应用《代数》第 IV 章第2节第3号的命题9看出这一点。因此
\end{enumerate}

\[
(\delta \circ \varphi) (u) = (\delta \circ \varphi) (u) \circ w,
\]

这证明了 \((\delta \circ \varphi)(\mathbf{Z}) \subset \mathbf{S}(\mathfrak{h})^{\mathrm{W}}\)。

\begin{enumerate}
\def\labelenumi{\alph{enumi})}
\setcounter{enumi}{1}
\tightlist
\item
  令 \(\eta\) 为第3号定理1的推论2中定义的从 \(\mathrm{I}(\mathfrak{g})\) 到 \(\mathbf{S}(\mathfrak{h})^{\mathrm{W}}\) 的同构。考虑从 g-模 U 到 g-模 \(\mathbf{S}(\mathfrak{g})\) 的典则同构（第 I 章第2节第8号），并令 \(\theta\) 为其在 Z 上的限制。则 \(\theta(\mathrm{Z}) = \mathrm{I}(\mathfrak{g})\) 。令 z 为 U 中滤子次数 \(\leq f\) 的 Z 元素。
\end{enumerate}

\[
\begin{array}{c c c} \mathrm{Z} & \stackrel {{\theta}} {{\longrightarrow}} & \mathrm{I} (\mathfrak {g}) \\ \varphi \Big \downarrow & & \Big \downarrow^ {\eta} \\ \mathbf {S} (\mathfrak {h}) & \stackrel {{\delta}} {{\longrightarrow}} & \mathbf {S} (\mathfrak {h}). \end{array}
\]

引入 §6 第4号的记号，并置

\[
z = \sum_ {\sum q _ {i} + \sum m _ {i} + \sum p _ {i} \leq f} \lambda_ {(q _ {i}), (m _ {i}), (p _ {i})} u ((q _ {i}), (m _ {i}), (p _ {i})).
\]

令 \(v((q_i),(m_i),(p_i))\) 为在 \(X_{-\alpha_1}^{q_1}\ldots X_{-\alpha_n}^{q_n}H_1^{m_1}\ldots H_l^{m_l}X_{\alpha_1}^{p_1}\ldots X_{\alpha_n}^{p_n}\) 中计算的单项式 \(\mathbf{S}(\mathfrak{g})\) 。记 \(\mathbf{S}_d(\mathfrak{g})\) 为 \(\mathbf{S}(\mathfrak{g})\) 中次数为 \(0,1,\ldots,d\) 的齐次分量之和，则有

\[
\theta (z) \equiv \sum_ {\sum q _ {i} + \sum m _ {i} + \sum p _ {i} = f} \lambda_ {(q _ {i}), (m _ {i}), (p _ {i})} v ((q _ {i}), (m _ {i}), (p _ {i})) \pmod {\mathbf {S} _ {f - 1} (\mathfrak {g})}
\]

因此

\[
(\eta \circ \theta) (z) \equiv \sum_ {\sum m _ {i} = f} \lambda_ {(0), (m _ {i}), (0)} v ((0), (m _ {i}), (0)) \pmod {\mathbf {S} _ {f - 1} (\mathfrak {h})}
\]

从而

\[
(\eta \circ \theta) (z) \equiv \varphi (z) \pmod {\mathbf {S} _ {f - 1} (\mathfrak {h})}.\tag{3}
\]

\begin{enumerate}
\def\labelenumi{\alph{enumi})}
\setcounter{enumi}{2}
\tightlist
\item
  我们证明 \(\delta \circ \varphi : \mathbf{Z} \to \mathbf{S}(\mathfrak{h})^{\mathrm{W}}\) 是双射。U 和 \(\mathbf{S}(\mathfrak{g})\) 上的典则滤子在 Z、\(\mathrm{I}(\mathfrak{g})\) 和 \(\mathbf{S}(\mathfrak{h})^{\mathrm{W}}\) 上诱导出滤子，而 \(\theta, \eta\) 与这些滤子相容，因而 \(\mathrm{gr}(\eta\circ\theta)\) 是从向量空间 \(\mathrm{gr}(\mathbf{Z})\) 到向量空间 \(\mathrm{gr}(\mathbf{S}(\mathfrak{h})^{\mathrm{W}})\) 的同构。由 (3)，\(\mathrm{gr}(\varphi)=\mathrm{gr}(\eta\circ\theta)\)，且显然 \(\mathrm{gr}(\delta)\) 是恒等映射。因此 \(\mathrm{gr}(\delta\circ\varphi)\) 是双射，所以
\end{enumerate}

\[
\delta \circ \varphi : Z \to \mathbf {S} (\mathfrak {h}) ^ {W}
\]

是双射（《交换代数》第 III 章第2节第8号定理1的推论1和2）。d）回忆 a）中的记号。令 E 为最高权为 \(\lambda\) 的单 g-模，\(\chi\) 为其中心特征标（§6，第1号，定义2）。令 \(\varphi'\) 和 \(\delta'\) 为相对于基 \(\varphi\) 而与 \(\delta\) 和 \(w(\mathrm{B})\) 类似的同态。E 相对于 \(w(\mathrm{B})\) 的最高权为 \(w(\lambda)\)。由 §6，第4号，命题7，

\[
\varphi (u) (\lambda) = \chi (u) = \varphi^ {\prime} (u) (w \lambda)
\]

对所有 \(u\in \mathbf{Z}\) 都成立，所以由 a）得，

\[
\begin{array}{c} (\delta \circ \varphi) (u) (w \lambda + w \rho) = (\delta \circ \varphi) (u) (\lambda + \rho) = \varphi (u) (\lambda) = \varphi^ {\prime} (u) (w \lambda) \\ = (\delta^ {\prime} \circ \varphi^ {\prime}) (u) (w \lambda + w \rho). \end{array}
\]

因此，多项式函数 \((\delta \circ \varphi)(u)\) 和 \((\delta' \circ \varphi')(u)\) 在 \(w(\mathrm{P}_{++}) + w\rho\) 上相同，因而相等。

推论 1。对任意 \(\lambda \in \mathfrak{h}^*\)，令 \(\chi_{\lambda}\) 为从 \(z \mapsto (\varphi(z))(\lambda)\) 到 \(Z\) 的同态 \(k\)。

\begin{enumerate}
\def\labelenumi{(\roman{enumi})}
\item
  若 \(k\) 是代数闭域，则从 \(Z\) 到 \(k\) 的每个同态都形如 \(\chi_{\lambda}\)，其中 \(\lambda \in \mathfrak{h}^*\)。
\item
  令 \(\lambda, \mu \in \mathfrak{h}^*\)。则 \(\chi_{\lambda} = \chi_{\mu}\) 当且仅当 \(\mu + \rho \in \mathrm{W}(\lambda + \rho)\)。
\end{enumerate}

若 \(k\) 是代数闭域，则从 \(\mathbf{S}(\mathfrak{h})^{\mathrm{W}}\) 到 \(k\) 的每个同态都可延拓为从 \(\mathbf{S}(\mathfrak{h})\) 到 \(k\) 的同态（《交换代数》第 V 章第1节第9号命题22，以及第2节第1号定理1的推论4），并且从 \(\mathbf{S}(\mathfrak{h})\) 到 \(k\) 的每个同态都形如 \(f \mapsto f(\lambda)\)，其中 \(\lambda \in \mathfrak{h}^*\)（第 VII 章附录 I 命题1）。因此，若 \(\chi\) 是从 \(Z\) 到 \(k\) 的同态，则存在（定理 2）某个 \(\mu \in \mathfrak{h}^*\)，使得对所有 \(z \in Z\)，

\[
\chi (z) = ((\delta \circ \varphi) (z)) (\mu) = (\varphi (z)) (\mu - \rho)
\]

因而得到 (i)。

令 \(\lambda, \mu \in \mathfrak{h}^*\)，并假设 \(\chi_{\lambda} = \chi_{\mu}\)。则对所有 \(z \in \mathbf{Z}\)，

\[
((\delta \circ \varphi) (z)) (\lambda + \rho) = (\varphi (z)) (\lambda) = \chi_ {\lambda} (z) = \chi_ {\mu} (z) = ((\delta \circ \varphi) (z)) (\mu + \rho);
\]

换言之，由 \(\mathbf{S}(\mathfrak{h})\) 和 \(\lambda + \rho\) 定义的从 \(\mu + \rho\) 到 k 的同态在 \(\mathbf{S}(\mathfrak{h})^{\mathrm{W}}\) 上相同；因此，断言 (ii) 由《交换代数》第 V 章第2节第2号定理2的推论得到。

推论 2。令 E、E' 为有限维单 g-模，\(\chi, \chi'\) 为它们的中心特征标。若 \(\chi = \chi'\)，则 E 与 E' 同构。

令 \(\lambda, \lambda'\) 为 E、E' 的最高权。由 §6，第4号，命题 7，\(\chi_{\lambda} = \chi = \chi' = \chi_{\lambda'}\)，所以存在 \(w \in \mathrm{W}\) 使 \(\lambda' + \rho = w(\lambda + \rho)\)。由于 \(\lambda + \rho\) 和 \(\lambda' + \rho\) 属于由 B 定义的房，故 \(w = 1\)。于是 \(\lambda = \lambda'\)，推论得证。

命题 7。对有限维单 g-模的任意类 \(\gamma\)，令 \(U_{\gamma}\) 为 g-模 U（取 g 在 U 上的伴随表示）中类型为 \(\gamma\) 的同构型分支。令 \(\gamma_{0}\) 为一维平凡 g-模的类。令 \([U,U]\) 为由 U 中元素对的括号生成的 U 的向量子空间。

\begin{enumerate}
\def\labelenumi{(\roman{enumi})}
\item
  U 是各 \(\mathrm{U}_{\gamma}\) 的直和。
\item
  \(U_{\gamma_{0}} = Z\)，且 \(\sum_{\gamma \neq \gamma_{0}} U_{\gamma} = [U, U]\)。
\item
  令 \(u \mapsto u^{\natural}\) 为由分解 \(\mathrm{U} = \mathrm{Z} \oplus [\mathrm{U}, \mathrm{U}]\) 定义的 U 到 Z 的投影。若 \(u \in \mathrm{U}\) 且 \(v \in \mathrm{U}\)，则 \((uv)^{\natural} = (vu)^{\natural}\)。若 \(u \in \mathrm{U}\) 且 \(z \in \mathrm{Z}\)，则 \((uz)^{\natural} = u^{\natural}z\)。
\item
  令 \(\varphi\) 为 Harish-Chandra 同态。令 \(\lambda \in \mathrm{P}_{++}\)，并令 E 为最高权为 \(\mathfrak{g}\) 的有限维单 \(\lambda\)-模。对所有 \(u \in \mathrm{U}\)，有
\end{enumerate}

\[
\frac {1}{\dim \mathrm{E}} \mathrm{Tr} (u _ {\mathrm{E}}) = (\varphi (u ^ {\natural})) (\lambda).
\]

 \(\mathfrak{g}\)-模 U 是有限维子模的直和。这蕴含 (i)。

显然 \(U_{\gamma_{0}} = Z\)。令 \(U'\) 为 U 的一个向量子空间，它在伴随表示下定义一个类为 \(\gamma\) 的子表示。于是或者 \([g, U'] = U'\)，或者 \([g, U'] = 0\)。因此，若 \(\gamma \neq \gamma_{0}\)，则 \([g, U'] = U'\)，从而 \(\sum_{\gamma \neq \gamma_{0}} U_{\gamma} \subset [U, U]\)。另一方面，若 \(u \in U\) 且 \(x_{1}, \ldots, x_{n} \in g\)，则

\[
\begin{array}{c} [ x _ {1} \ldots x _ {n}, u ] = (x _ {1} \ldots x _ {n} u - x _ {2} \ldots x _ {n} u x _ {1}) + (x _ {2} \ldots x _ {n} u x _ {1} - x _ {3} \ldots x _ {n} u x _ {1} x _ {2}) \\ + \dots + (x _ {n} u x _ {1} \ldots x _ {n - 1} - u x _ {1} \ldots x _ {n}) \in [ \mathfrak {g}, \mathrm{U} ]. \end{array}
\]

因此 \([\mathrm{U},\mathrm{U}]\subset \left[\mathfrak{g},\sum_{\gamma}\mathrm{U}_{\gamma}\right] = \left[\mathfrak{g},\sum_{\gamma \neq \gamma_0}\mathrm{U}_\gamma \right]\subset \sum_{\gamma \neq \gamma_0}\mathrm{U}_\gamma .\) 这证明了 (ii)。在这些条件下，(iii) 由第 I 章第6节第9号引理5得到。

最后，令 E、\(\lambda\) 如 (iv) 中所述。则

\[
\begin{array}{r l} & {\mathrm{Tr} (u _ {\mathrm{E}}) = \mathrm{Tr} ((u ^ {\natural}) _ {\mathrm{E}}) \qquad \text {since} u - u ^ {\natural} \in [ \mathrm{U}, \mathrm{U} ]} \\ & {\qquad = \mathrm{Tr} (\varphi (u ^ {\natural}) (\lambda). 1) \qquad (\S 6, \text {no.} 4, \text {Prop.} 7)} \\ & {\qquad = (\dim \mathrm{E}). \varphi (u ^ {\natural}) (\lambda).} \end{array}
\]

